\chapter{Intraday Machine-Learned Alphas}\label{ch:s1:intraday-machine-learned-alphas}

A gradient-boosted model trained on fourteen hours of a simulated order book predicts the next five seconds of the mid-price with an out-of-sample
R-squared of 25\%, the next thirty seconds with 5\%, and the next two minutes not at all. Traded naively (crossing the spread on every forecast of a fifth
of a tick or more) the five-second model loses 0.6 ticks a share: the spread it pays is larger than the moves it predicts. Traded passively, joining the
queue on the forecast's side and waiting to be filled, the same forecasts earn about a third of a tick a share. An \omterm{def:s1:intraday-machine-learned-alphas:alpha}{intraday alpha} is a prediction and an
execution rule together; neither is worth anything alone. The build is \texttt{firm.intraml}.

\section{Targets and horizons}

\begin{definition}[Intraday alpha, prediction horizon]\label{def:s1:intraday-machine-learned-alphas:alpha}
An \emph{intraday alpha}\index{intraday alpha} is a forecast of a security's price change over seconds to hours, built from market data (the order book, the
trades, related securities) and traded within the day. Its \emph{prediction horizon}\index{prediction horizon} is the interval over which the price change is
forecast, which sets how often the alpha trades and what execution cost it can bear.
\end{definition}

The target of this chapter is the change of the mid-price over the next 5, 30 and 120 seconds, in ticks, sampled every second on twenty one-hour sessions of
\texttt{firm.tape} (Book~7, chapter~2), whose default configuration includes a news window in mid-session. Its spread is about one tick (1.07 on average) and
the standard deviation of the target is 0.69 ticks at five seconds, 2.19 at thirty and 4.47 at two minutes. Kolm, Turiel and Westray, forecasting 115 Nasdaq
stocks at several horizons, found that stock-specific forecasts had an effective horizon of about two average price changes: beyond that, the book says
little. The synthetic tape agrees in shape.

\section{Features from the book and the tape}

Ten features, all computed from data up to the forecast's second (\cref{lst:s1:intraday-machine-learned-alphas:features}): the top-of-book imbalance, the
spread, the order-flow imbalance of Cont, Kukanov and Stoikov over the last 1, 5 and 30 seconds (Book~7, chapter~8), the signed trade volume over 5 and 30
seconds (Book~7, chapter~9), the mid's own change over 5 and 30 seconds, and the microprice's distance from the mid. Order-flow features, normalised by the
book's depth, are the ones the literature finds stationary: Kolm and co-authors report that models trained on order flow outperform most models trained on
raw book states. The unit tests check causality by construction: cutting the tape at a time leaves every feature before it unchanged.

\section{Models and their validation}

Two models per horizon: ridge regression on the standardised features, and a gradient-boosted ensemble (Friedman's method) of sixty depth-two regression
trees split on quantiles of the features, written in a few dozen lines of NumPy. Training uses sessions 1 to 14 and testing sessions 15 to 20: whole sessions
apart, so that overlapping targets (a thirty-second target at second $t$ shares 29 seconds with the one at $t + 1$) cannot leak across the split, the purging
and embargo of Book~7, chapter~20 at session scale.

\begin{center}
\small
\begin{tabular}{@{}lccc@{}}
\toprule
horizon & 5 seconds & 30 seconds & 120 seconds\\
\midrule
out-of-sample R-squared: ridge; boosted trees & 21.1\%; 24.6\% & 5.1\%; 5.1\% & $-0.4\%$; $-1.5\%$\\
rank IC: ridge; boosted trees & 0.33; 0.34 & 0.23; 0.28 & 0.07; 0.07\\
in-sample R-squared, boosted trees & 23.3\% & 9.2\% & 3.2\%\\
\midrule
every forecast above 0.2 ticks, aggressive (ticks a share) & $-0.60$ & $-0.48$ & $-0.87$\\
every forecast above 0.2 ticks, passive & 0.37 & 0.33 & $-0.02$\\
coupled: aggressive above the spread, passive below & 0.33 & 0.22 & $-0.24$\\
\bottomrule
\end{tabular}
\end{center}

The five-second R-squared of a quarter is high by the standard of real markets, and the reason is the tape: Book~7 found that its mids trend over a few
seconds, a property the chapter's model learns. What carries over is the pattern across horizons (a rapid decay of predictability, with the thirty-second
forecast already down to 5\% and the two-minute one worthless out of sample while it still fits 3.2\% in sample) and the small gain of the trees over ridge
where there is a signal to find. Sirignano and Cont found, with deep learning on billions of US quotes, that a model trained on all stocks predicts better
than stock-specific ones, even for stocks outside its training set: price formation from order flow is largely universal. That is an argument for pooling
data across instruments, which this chapter's twenty sessions only mimic.

\section{Coupling the alpha to execution}

\begin{definition}[Execution coupling]\label{def:s1:intraday-machine-learned-alphas:coupling}
\emph{Execution coupling}\index{execution coupling} is the design of an intraday strategy's orders from its forecast: crossing the spread only when the
forecast exceeds the cost of crossing, resting passively when it does not, and sizing and cancelling orders as the forecast changes, so that the alpha and the
execution are one decision.
\end{definition}

The table's lower half is the chapter's point. A forecast of a fifth of a tick does not pay a spread of a tick: crossing on every such forecast loses 0.60
ticks a share at five seconds. The same forecasts, used to choose which side of the book to join, earn 0.37 ticks a share on the orders that are filled
(\cref{lst:s1:intraday-machine-learned-alphas:passive}), because a passive order earns the spread instead of paying it and the forecast tilts the fills
toward the favourable side. Coupling the two (crossing only when the forecast exceeds the spread, 324 times over six test hours, and resting otherwise)
earns 0.33 ticks a share at five seconds and 0.22 at thirty. At two minutes nothing works: the forecast has no out-of-sample power, and both execution styles
lose. The passive numbers are optimistic in one way (the fill model fills a resting order as soon as the market trades through its price, ignoring its place
in the queue; Book~7, chapter~18 has the queue-aware replay) and honest in another (the exit crosses the spread).

\begin{omfigure}
\begin{tikzpicture}
\begin{axis}[name=a,ybar,bar width=8pt,width=6.4cm,height=5.2cm,xlabel={\small horizon (seconds)},ylabel={\small out-of-sample R-squared (\%)},
  tick label style={font=\footnotesize},symbolic x coords={5,30,120},xtick=data,ymin=-3,ymax=28,grid=major,grid style={gray!25},
  enlarge x limits=0.3,legend style={font=\scriptsize,at={(0.97,0.97)},anchor=north east},table/col sep=comma]
\addplot[fill=black!25,draw=black!60] table[x=h,y=ridge]{figdata/strategies-1/14-intraday-machine-learned-alphas/r2.csv};
\addplot[fill=omDef!60,draw=omDef] table[x=h,y=gbm]{figdata/strategies-1/14-intraday-machine-learned-alphas/r2.csv};
\legend{ridge,boosted trees}
\end{axis}
\begin{axis}[at={(a.east)},anchor=west,xshift=2.0cm,ybar,bar width=6pt,width=6.4cm,height=5.2cm,xlabel={\small horizon (seconds)},
  ylabel={\small P\&L per share (ticks)},tick label style={font=\footnotesize},symbolic x coords={5,30,120},xtick=data,ymin=-1,ymax=0.6,
  grid=major,grid style={gray!25},enlarge x limits=0.3,legend style={font=\scriptsize,at={(0.5,1.03)},anchor=south,legend columns=3},
  table/col sep=comma]
\addplot[fill=omThm!60,draw=omThm] table[x=h,y=naive_aggressive]{figdata/strategies-1/14-intraday-machine-learned-alphas/pnl.csv};
\addplot[fill=omDef!40,draw=omDef] table[x=h,y=passive]{figdata/strategies-1/14-intraday-machine-learned-alphas/pnl.csv};
\addplot[fill=omDef!80,draw=omDef] table[x=h,y=coupled]{figdata/strategies-1/14-intraday-machine-learned-alphas/pnl.csv};
\legend{aggressive,passive,coupled}
\end{axis}
\end{tikzpicture}

\omcaption{Machine-learned forecasts of the synthetic tape's mid-price. Left: out-of-sample R-squared by horizon. Right: P\&L per share of the boosted-tree
forecasts on six test hours, traded aggressively, passively, or coupled (aggressive only above the spread). Data: \texttt{s1\_intraml}.}
\label{fig:s1:intraday-machine-learned-alphas:results}
\end{omfigure}

\section{Strategy files}

\begin{strategyfile}{Book-imbalance alpha}\label{strat:s1:intraday-machine-learned-alphas:book}
\sfield{Who pays you, and why} Traders whose orders reveal pressure before prices move: the book's imbalance is the queue that will be eaten next.
\sfield{Instruments and venues} Liquid stocks and futures with deep, visible books.
\sfield{Signal} Top-of-book and depth imbalance, order-flow imbalance over seconds.
\sfield{Sizing and execution} Mostly passive, joining the favoured side; aggressive only when the forecast beats the spread.
\sfield{Costs} Spread, fees and rebates; queue position.
\sfield{How it dies} Faster competitors reading the same book; spoofed depth.
\sfield{Horizon, capacity, infrastructure} Seconds; small capacity per name; co-located, low-latency systems.
\sfield{Backtest honestly} A queue-aware replay (Book~7, chapter~18); latency; fees.
\sfield{Sources} Cont, Kukanov and Stoikov (2014); this chapter's simulation.
\end{strategyfile}

\begin{strategyfile}{Trade-flow alpha}\label{strat:s1:intraday-machine-learned-alphas:flow}
\sfield{Who pays you, and why} Metaorders split into many trades, whose signs persist (Book~7, chapter~9).
\sfield{Instruments and venues} As for book imbalance.
\sfield{Signal} Signed trade volume over seconds to minutes; the persistence of trade signs.
\sfield{Sizing and execution} As for book imbalance.
\sfield{Costs} As for book imbalance.
\sfield{How it dies} Metaorders that hide better; competitors on the same flow.
\sfield{Horizon, capacity, infrastructure} Seconds to minutes; the trade feed with its sign.
\sfield{Backtest honestly} Trade signs as they could be known (Lee--Ready or the feed's aggressor flag), with its latency.
\sfield{Sources} Kolm, Turiel and Westray (2023) on order-flow inputs.
\end{strategyfile}

\begin{strategyfile}{Cross-asset lead alpha}\label{strat:s1:intraday-machine-learned-alphas:lead}
\sfield{Who pays you, and why} Slower instruments that follow faster ones (a stock after its index future, an ETF after its constituents).
\sfield{Instruments and venues} Pairs of related instruments on the same or different venues.
\sfield{Signal} The leader's recent move, relative to the follower's (Book~7, chapter~10).
\sfield{Sizing and execution} Trade the follower in the leader's direction before it catches up.
\sfield{Costs} Speed: the gain is only there for whoever is first.
\sfield{How it dies} Latency competition; the lead shrinking to microseconds.
\sfield{Horizon, capacity, infrastructure} Milliseconds to seconds; networks and co-location (Books~10 and~11).
\sfield{Backtest honestly} Timestamps from the same clock; realistic latency between venues.
\sfield{Sources} Sirignano and Cont (2019) on universal, pooled features; no performance figure verified.
\end{strategyfile}

\begin{strategyfile}{Alpha-driven execution}\label{strat:s1:intraday-machine-learned-alphas:execution}
\sfield{Who pays you, and why} The spread, earned instead of paid, when the forecast chooses where to rest.
\sfield{Instruments and venues} Any instrument the firm trades for other reasons, and market-making books.
\sfield{Signal} The short-horizon forecast, against the spread and the queue.
\sfield{Sizing and execution} Aggressive above the spread, passive below, cancel when the forecast turns.
\sfield{Costs} Adverse selection on passive fills; cancellation limits.
\sfield{How it dies} It does not die; it saves cost on every other strategy's trading.
\sfield{Horizon, capacity, infrastructure} Seconds; an execution engine that takes forecasts as input.
\sfield{Backtest honestly} Queue-aware fills; the markouts of passive fills (Book~7, chapter~23).
\sfield{Sources} This chapter's simulation (0.33 ticks a share coupled against $-0.60$ aggressive at five seconds).
\end{strategyfile}

\section{Tutorial: one per cent is a lot}\label{tut:s1:intraday-machine-learned-alphas:tut}

\begin{tutorial}
\textbf{Goal.} Build features from the synthetic tape, train and validate two models at three horizons, and trade their forecasts three ways. \textbf{End state:}
the table and \cref{fig:s1:intraday-machine-learned-alphas:results}.
\begin{tutsteps}
\item \textbf{Features}: every second, from data up to that second.
\omcode{code/firm/intraml/firm_intraml.py}{34}{59}{Order-book and trade-flow features.}\label{lst:s1:intraday-machine-learned-alphas:features}
\item \textbf{Passive execution}: join the touch on the forecast's side; filled only if traded through; exit across the spread.
\omcode{code/firm/intraml/firm_intraml.py}{154}{172}{Passive entries with a touch-fill model.}\label{lst:s1:intraday-machine-learned-alphas:passive}
\item \textbf{Run} \texttt{scores(h)}, \texttt{naive(h)} and \texttt{trading(h)} for 5, 30 and 120 seconds, and \texttt{fig\_intraml.py}.
\end{tutsteps}
\textbf{What to change next.} Replace the touch-fill model by Book~7's queue-aware replay; add the second instrument of \texttt{simulate\_pair} as a lead
feature; train one model on all horizons with the horizon as a feature.
\end{tutorial}

\section{Build: intraday machine learning}\label{bld:s1:intraday-machine-learned-alphas:intraml}

\begin{build}
\bfield{Purpose} Causal intraday features, horizon targets, two models with session-level validation, and execution rules coupled to the forecasts.
\bfield{Interface} \texttt{features(tape, step)}, \texttt{targets(tape, times, horizons)}, \texttt{ridge} and \texttt{ridge\_predict}, \texttt{gbm} and
\texttt{gbm\_predict}, \texttt{r2(y, yhat)}, \texttt{aggressive(pred, fwd, spread, threshold)}, \texttt{passive(tape, times, pred, h, threshold, wait)}.
\bfield{Rules} Features from data up to each time; training and test separated by whole sessions; costs in every P\&L.
\bfield{Acceptance tests} \texttt{code/firm/intraml/tests/}: features unchanged when the future is removed; both models recover a planted signal; the aggressive
P\&L by hand.
\bfield{Stretch} A queue-aware passive fill model; cross-asset features; deeper trees and early stopping on a validation session.
\end{build}

\begin{omsources}
\item R.~Cont, A.~Kukanov and S.~Stoikov, ``The price impact of order book events'', \emph{Journal of Financial Econometrics} 12(1), 2014.
\item P.~N.~Kolm, J.~Turiel and N.~Westray, ``Deep order flow imbalance'', \emph{Mathematical Finance} 33(4), 2023.
\item J.~Sirignano and R.~Cont, ``Universal features of price formation in financial markets'', \emph{Quantitative Finance} 19(9), 2019.
\item J.~H.~Friedman, ``Greedy function approximation: a gradient boosting machine'', \emph{Annals of Statistics} 29(5), 2001.
\end{omsources}

\section{Exercises}

\begin{exercise}[$\star$]\label{exo:s1:intraday-machine-learned-alphas:1}
A forecast has an R-squared of 25\% against a target with a standard deviation of 0.69 ticks. What is the standard deviation of the forecast?
\end{exercise}

\begin{exercise}[$\star$]\label{exo:s1:intraday-machine-learned-alphas:2}
Why can a thirty-second target at second $t$ not be validated on a sample that contains second $t + 1$?
\end{exercise}

\begin{exercise}[$\star$]\label{exo:s1:intraday-machine-learned-alphas:3}
An aggressive round trip pays one spread of a tick. What must a forecast predict, in ticks, to break even?
\end{exercise}

\begin{exercise}[$\star\star$]\label{exo:s1:intraday-machine-learned-alphas:4}
Why do passive fills earn money on forecasts too small to pay the spread? What biases the chapter's passive numbers upward?
\end{exercise}

\begin{exercise}[$\star\star$]\label{exo:s1:intraday-machine-learned-alphas:5}
The two-minute model fits 3.2\% in sample and $-1.5\%$ out of sample. What happened?
\end{exercise}

\begin{exercise}[$\star\star$]\label{exo:s1:intraday-machine-learned-alphas:6}
Why normalise order-flow features by the book's depth?
\end{exercise}

\begin{exercise}[$\star\star\star$]\label{exo:s1:intraday-machine-learned-alphas:7}
\emph{Coding.} Train the models with the mid's own past changes removed from the features, and compare the five-second R-squared. What does the difference
say about the synthetic tape?
\end{exercise}

\begin{exercise}[$\star\star\star$]\label{exo:s1:intraday-machine-learned-alphas:8}
\emph{Find the flaw.} ``Our model predicts the next second's price change with 60\% accuracy on direction; we will trade it with market orders on 3,000
stocks.''
\end{exercise}

\section{Problem: One Per Cent Is a Lot}

\begin{problem}[{Weekend problem --- a forecast and its execution}]\label{pb:s1:intraday-machine-learned-alphas:1}
The chapter's models on \texttt{firm.tape}, and the published record.

\medskip\noindent\textbf{Part I --- The forecast.}
\begin{enumerate}
\item Define an \omterm{def:s1:intraday-machine-learned-alphas:alpha}{intraday alpha} and its \omterm{def:s1:intraday-machine-learned-alphas:alpha}{prediction horizon}.
\item Describe the targets, the sessions and the spread.
\item List the features and how causality is tested.
\item What did Kolm, Turiel and Westray find about horizons and inputs?
\end{enumerate}

\medskip\noindent\textbf{Part II --- The models.}
\begin{enumerate}[resume]
\item Describe the two models.
\item How are training and test separated, and why by session?
\item Give the out-of-sample R-squared and IC by horizon.
\item Why is the five-second R-squared so high here?
\end{enumerate}

\medskip\noindent\textbf{Part III --- Execution.}
\begin{enumerate}[resume]
\item Define \omterm{def:s1:intraday-machine-learned-alphas:coupling}{execution coupling}.
\item Give the P\&L per share of the three execution styles by horizon.
\item Why do aggressive trades lose on small forecasts?
\item What does the passive fill model assume?
\end{enumerate}

\medskip\noindent\textbf{Part IV --- The verdict.}
\begin{enumerate}[resume]
\item State the \emph{named result}: out-of-sample R-squared by horizon and the P\&L per share after the spread, aggressive against passive.
\item What did Sirignano and Cont find about universality?
\item How would you pool data across instruments?
\item What would change with a queue-aware replay?
\item How does latency enter?
\item Which strategy file is least exposed to competition?
\item What is the capacity of a five-second alpha?
\item In one sentence: what is an \omterm{def:s1:intraday-machine-learned-alphas:alpha}{intraday alpha} worth?
\end{enumerate}
\end{problem}

\section{Interview questions}

\begin{interviewq}[$\star$ \iqroles{researcher, mle}]\label{iq:s1:intraday-machine-learned-alphas:1}
What features would you build from an order book to predict the next few seconds?
\end{interviewq}

\begin{interviewq}[$\star\star$ \iqroles{researcher, mle}]\label{iq:s1:intraday-machine-learned-alphas:2}
How do you cross-validate a model whose targets overlap in time?
\end{interviewq}

\begin{interviewq}[$\star\star$ \iqroles{trader}]\label{iq:s1:intraday-machine-learned-alphas:3}
Your model predicts moves smaller than the spread. Is it useless?
\end{interviewq}

\begin{interviewq}[$\star\star$ \iqroles{mle, developer}]\label{iq:s1:intraday-machine-learned-alphas:4}
How would you serve a gradient-boosted model within microseconds of a market data update?
\end{interviewq}

\begin{interviewq}[$\star\star$ \iqroles{researcher}]\label{iq:s1:intraday-machine-learned-alphas:5}
Why might an R-squared of 1\% be a good intraday model?
\end{interviewq}

\begin{interviewq}[$\star\star\star$ \iqroles{researcher}]\label{iq:s1:intraday-machine-learned-alphas:6}
A forecast $f$ of a move $y$ has correlation $\rho$ with it and both are normal with standard deviations $\sigma_f$ and $\sigma_y$. With a round-trip cost
$c$, derive the expected P\&L per trade of trading only when $|f| > c$, and discuss how it depends on $\rho$.
\end{interviewq}
